\chapter{Implementation, Experiments, and Evaluation}
\label{chap:experiments}

\section{Implementation}

The main experiment runner is implemented in \texttt{src/kge\_nfm.py}. For each
fold, the runner loads data, trains or prepares the KGE model, extracts pair
features, trains the NFM classifier, evaluates predictions, and writes output
artifacts. The project is implemented as a Python-based machine learning
pipeline using PyTorch for supervised learning and PyKEEN for knowledge graph
embedding models.

\subsection{Environment and Libraries}

\begin{table}[H]
\centering
\caption{Implementation environment}
\label{tab:development_environment}
\begin{tabularx}{\textwidth}{lX}
\toprule
\textbf{Item} & \textbf{Description} \\
\midrule
Language & Python 3.10 or newer \\
Deep learning framework & PyTorch with the CUDA 12.1 (\texttt{cu121}) wheel build \\
KGE framework & PyKEEN \\
Data processing & NumPy, pandas, and scikit-learn \\
Experiment output & CSV metrics, predictions, curves, and model checkpoints \\
\bottomrule
\end{tabularx}
\end{table}

PyTorch provides tensor operations, neural network training, and checkpoint
management. PyKEEN is used to train the DistMult and CompGCN encoders. NumPy
and pandas support numerical and tabular data processing, while scikit-learn
provides preprocessing and evaluation utilities.

\subsection{Source-Code Organization}

\begin{table}[H]
\centering
\caption{Main implementation files}
\label{tab:source_files}
\begin{tabularx}{\textwidth}{lX}
\toprule
\textbf{Path} & \textbf{Purpose} \\
\midrule
\texttt{src/kge\_nfm.py} & Main KGE-NFM experiment runner \\
\texttt{src/data\_utils.py} & Dataset, fold, knowledge graph, and descriptor loading \\
\texttt{src/kge.py} & KGE training, scoring, and embedding extraction \\
\texttt{src/nfm.py} & NFM architecture and training loop \\
\texttt{src/metrics\_utils.py} & ROC-AUC and PR-AUC computation \\
\texttt{src/utils.py} & Random seed, device, and output-directory utilities \\
\texttt{baseline/} & Baseline implementations and exploratory experiments \\
\texttt{output/} & Predictions, metrics, curves, and saved models \\
\bottomrule
\end{tabularx}
\end{table}

\subsection{Experiment Workflow}

The experiment workflow is:

\begin{enumerate}
    \item parse command-line arguments;
    \item set random seed and select computation device;
    \item load dataset specification and fold files;
    \item load descriptor features and knowledge graph triples;
    \item train the selected KGE model;
    \item compute KGE scores and entity embeddings;
    \item construct train and test feature matrices;
    \item train the NFM classifier;
    \item compute ROC-AUC and PR-AUC;
    \item save predictions, metric curves, and model checkpoints.
\end{enumerate}

\subsection{Execution Command}

The following command runs the KGE-NFM pipeline with the CompGCN encoder on
the Yamanishi08 warm-start 1:10 setting. Replacing \texttt{compgcn} with
\texttt{distmult} selects the alternative encoder under the same downstream
configuration.

\begin{lstlisting}[caption={Example KGE-NFM experiment command}]
python src/kge_nfm.py \
  --dataset yamanishi_08 \
  --split warm_start_1_10 \
  --folds 10 \
  --kge-model compgcn \
  --device auto \
  --embedding-dim 400 \
  --kge-epochs 50 \
  --nfm-epochs 2000 \
  --output-dir output/kge_nfm_torch
\end{lstlisting}

\section{Evaluation Metrics}

\subsection{ROC-AUC}

ROC-AUC measures how well the model ranks positive interactions above negative
or unknown interactions across all classification thresholds.

\subsection{PR-AUC}

PR-AUC summarizes the precision-recall curve. This metric is especially useful
for DTI prediction because positive interactions are usually sparse.

\section{Baselines}

% TODO: Add the actual baselines evaluated in this project.

\begin{table}[H]
\centering
\caption{Baseline models}
\label{tab:baseline_models}
\begin{tabularx}{\textwidth}{lX}
\toprule
\textbf{Model} & \textbf{Description} \\
\midrule
Descriptor-only classifier & Uses drug and protein descriptors without KGE features \\
Standalone KGE score & Uses only KGE interaction scores for prediction \\
KGE + Random Forest & Uses KGE features with a random forest classifier \\
KGE + XGBoost & Uses KGE features with an XGBoost classifier \\
KGE-NFM & Framework combining KG embeddings, descriptors, and NFM \\
\bottomrule
\end{tabularx}
\end{table}

\section{Experimental Setup}

\begin{table}[H]
\centering
\caption{Experiment setup}
\label{tab:experiment_setup}
\begin{tabularx}{\textwidth}{lX}
\toprule
\textbf{Item} & \textbf{Setting} \\
\midrule
Dataset & Yamanishi08 \\
Splits & \texttt{warm\_start\_1\_1} and \texttt{warm\_start\_1\_10} \\
Number of supplied splits & 10 per setting \\
KGE encoders & DistMult and CompGCN \\
Random seed & 0 \\
Device & Automatically selected CPU or CUDA device \\
Output directory & \texttt{output/kge\_nfm\_torch} \\
\bottomrule
\end{tabularx}
\end{table}

\section{Main Results}

% TODO: Fill this table using the final output CSV files.

\begin{table}[H]
\centering
\caption{Main DTI prediction results}
\label{tab:main_results}
\begin{tabular}{lcccc}
\toprule
\textbf{Model} & \textbf{ROC-AUC} & \textbf{PR-AUC} &
\textbf{ROC-AUC Std.} & \textbf{PR-AUC Std.} \\
\midrule
Descriptor-only classifier & -- & -- & -- & -- \\
Standalone KGE & -- & -- & -- & -- \\
KGE + Random Forest & -- & -- & -- & -- \\
KGE + XGBoost & -- & -- & -- & -- \\
KGE-NFM & -- & -- & -- & -- \\
\bottomrule
\end{tabular}
\end{table}

\section{Ablation Study}

An ablation study helps measure the contribution of each feature source or
model component.

\begin{table}[H]
\centering
\caption{Ablation study template}
\label{tab:ablation_study}
\begin{tabular}{lcc}
\toprule
\textbf{Configuration} & \textbf{ROC-AUC} & \textbf{PR-AUC} \\
\midrule
Descriptors only & -- & -- \\
KGE embeddings only & -- & -- \\
Descriptors + KGE embeddings & -- & -- \\
Full KGE-NFM pipeline & -- & -- \\
\bottomrule
\end{tabular}
\end{table}

\section{Discussion}

% TODO: Discuss the final results. Suggested points:
% - Does KGE improve over descriptor-only baselines?
% - Does NFM improve over classical classifiers?
% - Are ROC-AUC and PR-AUC consistent?
% - Which split is hardest and why?

\section{Error Analysis}

% TODO: Add analysis of false positives and false negatives. Consider whether
% errors are concentrated around rare drugs, rare proteins, missing KG entities,
% or cold-start cases.

\section{Limitations}

The current project has several expected limitations:

\begin{itemize}
    \item negative labels may include unknown positive interactions;
    \item knowledge graph coverage may be incomplete for some entities;
    \item cold-start prediction is more difficult than warm-start prediction;
    \item hyperparameter search may be limited by available computation;
    \item model predictions require biological validation before practical use.
\end{itemize}
